\documentclass[12pt,a4paper]{article}
\usepackage[utf8]{inputenc}
\usepackage[french]{babel}
\usepackage[T1]{fontenc}
\usepackage{amsmath, amsfonts, amssymb}
\usepackage{graphicx}
\usepackage{geometry}
\usepackage{xcolor}
\usepackage{titlesec}
\usepackage{hyperref}
\usepackage{listings}
\usepackage{booktabs}
\usepackage{fancyhdr}

\geometry{margin=2.5cm}

\definecolor{codeblue}{RGB}{0,70,160}
\definecolor{codegray}{RGB}{245,245,245}

\lstset{
    basicstyle=\ttfamily\small,
    breaklines=true,
    frame=single,
    backgroundcolor=\color{codegray},
    rulecolor=\color{gray!50},
    keywordstyle=\bfseries\color{codeblue},
    commentstyle=\itshape\color{gray},
    showstringspaces=false,
    numbers=left,
    numberstyle=\tiny\color{gray},
    numbersep=8pt
}

\pagestyle{fancy}
\fancyhf{}
\rhead{NEUSSI Patrice -- 24P820}
\lhead{Realite Augmentee -- TP3}
\rfoot{\thepage}

\titleformat{\section}{\Large\bfseries\color{codeblue}}{}{0em}{\thesection\quad}
\titleformat{\subsection}{\large\bfseries}{}{0em}{\thesubsection\quad}
\hypersetup{colorlinks=true, linkcolor=codeblue, urlcolor=codeblue}

\begin{document}

\begin{titlepage}
    \centering
    \vspace*{2cm}
    {\huge\bfseries Rapport de Travaux Pratiques\par}
    \vspace{1cm}
    {\LARGE\bfseries\color{codeblue} TP3 -- Matrices 4x4 et Rasteriseur\par}
    \vspace{1.5cm}
    \rule{\linewidth}{0.5pt}\vspace{0.5cm}
    {\large Cours : Realite Augmentee FROM SCRATCH\par}
    {\large Module 03 -- NkMath\par}
    \vspace{0.5cm}\rule{\linewidth}{0.5pt}
    \vspace{2cm}
    {\large\textbf{Etudiant :} NEUSSI NJIETCHEU PATRICE EUGENE\par}
    {\large\textbf{Matricule :} 24P820\par}
    \vspace{2cm}
    {\large Repository GitHub :\par}
    \href{https://github.com/neussi/RA_TP3_Matrices_et_Rendu}{https://github.com/neussi/RA\_TP3\_Matrices\_et\_Rendu}
    \vfill
    {\large \today\par}
\end{titlepage}

\tableofcontents
\newpage

\section{Introduction}

Ce rapport couvre le module \textbf{Matrices 4x4 et Rasteriseur Logiciel} du cours
\textit{Realite Augmentee FROM SCRATCH}. Les matrices sont au coeur du pipeline graphique :
transformations, vues, projections.

\section{Concepts Theoriques}

\subsection{Matrice 4x4 (Mat4d)}

Une matrice $4\times4$ repr\'esente n'importe quelle transformation affine en homog\`ene.
Le produit matriciel :
\[
    C_{ij} = \sum_{k=0}^{3} A_{ik} \cdot B_{kj}
\]

\subsection{Transformations Fondamentales}

\begin{itemize}
    \item \textbf{Translation} : $T(t_x, t_y, t_z)$ -- colonne 3 de l'identit\'e
    \item \textbf{Rotation} autour d'un axe (formule de Rodrigues) :
    \[
        R = I\cos\theta + (1-\cos\theta)\hat{n}\hat{n}^T + \sin\theta[\hat{n}]_\times
    \]
    \item \textbf{Scale} : matrice diagonale $(s_x, s_y, s_z, 1)$
    \item \textbf{LookAt} : matrice de vue
    \item \textbf{Perspective} : projection perspective avec frustum
\end{itemize}

\subsection{Inverse d'une Matrice}

Methode de Gauss-Jordan sur la matrice augmentee $[M | I]$ :
\[
    [M | I] \longrightarrow [I | M^{-1}]
\]
Si le pivot est nul, la matrice est singuliere et \texttt{Inverse} retourne \texttt{false}.

\section{Realisation Technique}

\begin{lstlisting}[language=C++, caption=Produit matriciel Mat4d]
Mat4d operator*(const Mat4d& b) const {
    Mat4d r;
    for (int i = 0; i < 4; i++)
        for (int j = 0; j < 4; j++) {
            r(i,j) = 0;
            for (int k = 0; k < 4; k++)
                r(i,j) += (*this)(i,k) * b(k,j);
        }
    return r;
}
\end{lstlisting}

\begin{center}
\begin{tabular}{ll}
    \toprule
    \textbf{Fonction} & \textbf{Description} \\
    \midrule
    \texttt{Mat4d::Identity()} & Matrice identite \\
    \texttt{Mat4d::RotateAxis(n, theta)} & Rotation autour de l'axe n \\
    \texttt{Inverse(M, inv)} & Inverse par Gauss-Jordan \\
    \texttt{LookAt(eye, target, up)} & Matrice de vue \\
    \texttt{Perspective(fov, aspect, n, f)} & Projection perspective \\
    \texttt{TRS(T, R, S)} & Matrice TRS composee \\
    \texttt{DecomposeTRS(M, T, R, S)} & Decomposition \\
    \texttt{ProjectToScreen(p, w, h)} & Projection ecran \\
    \bottomrule
\end{tabular}
\end{center}

\section{Tests Unitaires (Framework Nkentseu)}

\subsection{Cas de test implementes}

\begin{itemize}
    \item \textbf{TP3\_IdentityMultiplication} : $M \times I = M$ pour 10 matrices aleatoires
    \item \textbf{TP3\_Inverse} : $M \times M^{-1} = I$ pour 10 matrices aleatoires
    \item \textbf{TP3\_SingularMatrix} : retourne false pour une matrice singuliere
    \item \textbf{TP3\_RotationAxis} : rotation de $\pi/2$ autour de Y de $(1,0,0) \to (0,0,-1)$
\end{itemize}

\subsection{Resultats d'execution}

\begin{lstlisting}[caption=Sortie de jenga test -- TP3]
Running test suite: TP3_Tests
Session started: 2026-04-05 16:52:46
Number of tests: 4

[OK]  TP3_IdentityMultiplication  10/10 assertions  (< 1ms)
[OK]  TP3_Inverse                 10/10 assertions  (< 1ms)
[OK]  TP3_SingularMatrix           1/1  assertions  (< 1ms)
[OK]  TP3_RotationAxis             3/3  assertions  (< 1ms)

-------------------------------------------------------
RESULTATS DES TESTS
  SUCCES
  Tests :       4 reussis, 4 au total
  Assertions :  24 reussies, 24 au total
  Taux succes : Tests: 100.0%, Assertions: 100.0%
  Temps total : 5ms (1ms/test)

Tous les tests sont reussis !
\end{lstlisting}

\section{Code Source}

\begin{center}
\href{https://github.com/neussi/RA_TP3_Matrices_et_Rendu}{https://github.com/neussi/RA\_TP3\_Matrices\_et\_Rendu}
\end{center}

\section{Conclusion}

Ce TP constitute l'epine dorsale du pipeline graphique. L'inverse par Gauss-Jordan et les
matrices de transformation (LookAt, Perspective) permettent de construire un rasteriseur
logiciel complet servant de base au TP suivant.

\end{document}
